\addchap{PREFACE}
The REXUS/BEXUS programme is realised under a bilateral Agency Agreement between the German Aerospace Center (DLR) and the Swedish National Space Agency (SNSA). The Swedish share of the payload has been made available to students from other European countries through a collaboration with the European Space Agency (ESA). EuroLaunch, a cooperation between the Swedish Space Corporation (SSC) and the Mobile Rocket Base (MORABA) of DLR, is responsible for the campaign management and operations of the launch vehicles. Experts from DLR, SSC, ZARM and ESA provide technical support to the student teams throughout the project. REXUS and BEXUS are launched from SSC, Esrange Space Center in northern Sweden.

\gls{mirage} is an experiment developed by a team of students from the Uppsala Universitet in Sweden. Its aim is to demonstrate low-cost methane sensing capabilities using stratospheric balloons. \gls{mirage} is scheduled to fly aboard BEXUS as part of cycle 17 of the Swedish-German programme REXUS/BEXUS.

This document contains all relevant available information regarding \gls{mirage}; including the experiment's objectives and requirements, a breakdown of the experimental module's design, design verification procedures, a timeline and pre-launch campaign activities. Technical documents that pertain to the design are also included for reference, as appendix sections.